\textit{To embrace fugitivity and resistance by and through technology becomes a unique form of media mediation. It establishes many ways of doing/performing what is known to be outside the margins of "normal" society.}

\subsection{Fugitivity as Technological Practice}(\parencite{Harney_Moten_2013}, \parencite{Moten_2003})
"Fugitivity" can be viewed as a methodology for working within and against institutionalized systems (an technologies). This work required "study" as an alternative to education, providing the needed understanding of how the marginalized create and communicate knowledge.

Bridge:

 If we buy into the media narrative tropes of those who operate with fugitive mindsets, they are resisting "The Man." However, whomever "Man" is has been studied and defined through scholarship about the dominant, institutional, paradigm.
 
\subsection{Sociogenic Critique}(\parencite{wynter2003}, \parencite{Wynter_2005})
There exists a critique of "Man," as the persistent over representation of an Eurocentric, Western "bourgeois." The systems and technologies that are championed by them are often presented as neutral, but instead encode their specific ways of knowing and being. Thus, the "others' engage in "sociogeny", their own cultural practices that create different way of being human, different from Man.

Bridge:

As the in-power have so often co-opted or appropriated those aspects of the marginalized knowledge-base that is most tasteful or desirable; so too must the marginalized appropriate, tactfully, the systems and technologies that render them invisible, in order to become undeniably clear and present.

\subsection{Black Radical Imagination and Technological Appropriation}(\parencite{Robinson_2020}, \parencite{Kelley_Monet_2003})
There already exists many historical precedencies for the tactical appropriation and transformation of dominant systems by Black communities. By examining them, and empowering them through a "Black Radical Imagination," we can establish alternative organizational and technological principles.

Bridge:

 This emergence of imagination alongside technological manipulation provides the foundation for engaging in a novel form of computational resistance.
 
\subsection{Digital Marronage and Computational Resistance}(\parencite{Johnson_2018})
Through "digital marronage," Black communities are already creating spaces of technological resistances through alternative practices. However, we must remain critical of the practices that "mark up" Black bodies and examine novel ways of understanding marginalized digital humanities.

Bridge to Next Section:

The body, as the site for knowledge and cultural memory must nor be overly marked or taxidermized, which renders static and immutable the materials of cultural knowledge and memory. Instead, we must mark-up and annotate the technologies, understanding their constraints and modifying them to better suit the needs of  the community.